\documentclass[11pt]{article}
\usepackage[a4paper,margin=22mm]{geometry}
\usepackage{fontspec}
\setmainfont{DejaVu Serif}
\setsansfont{DejaVu Sans}
\usepackage{microtype}
\usepackage{xcolor}
\usepackage{hyperref}
\usepackage{enumitem}
\usepackage{parskip}
\definecolor{Accent}{HTML}{971D32}
\definecolor{Muted}{HTML}{666666}
\hypersetup{colorlinks=true,urlcolor=Accent,linkcolor=Accent}
\setlist{leftmargin=1.6em,itemsep=0.3em,topsep=0.35em}
\setcounter{tocdepth}{2}
\renewcommand{\contentsname}{Contents}
\newcommand{\sectionrule}{\par\vspace{-0.5em}\noindent\textcolor{Accent}{\rule{\linewidth}{0.6pt}}\par}
\begin{document}

\begin{center}
{\sffamily\bfseries\color{Accent} TOEFL WORD OF THE DAY\par}
\vspace{8mm}
{\fontsize{42}{48}\selectfont\bfseries delegate\par}
\vspace{2mm}
{\Large\sffamily\color{Accent} /ˈdɛləɡət/\par}
\vspace{2mm}
{\small\sffamily Local audio phrase: a delegate\par}
{\small\sffamily Audio file: audios/a-delegate.mp3\par}
\vspace{3mm}
{\bfseries\color{Accent} Pronunciation changes across meanings.\par}
countable noun: /ˈdɛləɡət/ --- “a delegate”\par
transitive verb: /ˈdɛləˌɡeɪt/ --- “delegate tasks”\par
\vspace{2mm}
{\large\sffamily\color{Muted} countable noun and transitive verb\par}
\vspace{5mm}
{\sffamily a chosen representative \textbullet\ give someone a task \textbullet\ transfer authority\par}
\vspace{6mm}
{\large A delegate represents other people, while the verb delegate means giving another person a task or some authority.\par}
\vspace{6mm}
\begin{minipage}{0.84\textwidth}
\itshape “Effective managers know when to delegate routine tasks to other team members.”
\end{minipage}\par
\vspace{10mm}
\textbf{Date:} August 5th 2026\quad\textbullet\quad \textbf{Level:} B1--mid-B2 TOEFL
\vfill
Created and compiled by \textbf{Amir Kooshky}\par
\vspace{2mm}
\href{https://t.me/KooshkyTOEFL}{Telegram Channel}\quad\textbullet\quad
\href{https://www.instagram.com/kooshkytoefl}{Instagram}\quad\textbullet\quad
\href{https://t.me/KooshkyTOEFL\_pv}{Personal Telegram}
\end{center}
\newpage
\tableofcontents
\newpage

\section{Meanings and usage}
\sectionrule
\subsection{Meaning 1: Chosen representative}
\textbf{Part of speech:} countable noun

\textbf{IPA:} /ˈdɛləɡət/

\textbf{Pronunciation phrase:} a delegate

\textbf{Local audio file:} audios/a-delegate.mp3

\textbf{Register:} neutral to formal; common in politics, conferences, unions, institutions, and international organizations

\textbf{Definition:} a person chosen or elected to represent a group at a meeting, conference, or organization

\textbf{In simpler words:} a person who represents a group

\textbf{Typical pattern:} a delegate to + event or a delegate from + group or place

\subsubsection{Natural examples}
\begin{enumerate}
  \item Each country sent several delegates to the international conference.
  \item The union delegates voted in favor of the proposed agreement.
  \item She attended the convention as a delegate from her university.
  \item Delegates from more than fifty nations participated in the climate summit.
  \item The committee elected two student delegates to present its concerns.
  \item Conference delegates received copies of the final report.
  \item The delegate spoke on behalf of workers in the manufacturing sector.
\end{enumerate}

\subsubsection{Nuance and implication}
\textbf{Central nuance:} A delegate does not attend only as a private individual; the person represents the interests or views of a larger group.

\textbf{Usual implication:} The delegate may be expected to speak, negotiate, or vote on behalf of the group.

\subsubsection{Grammar notes}
\textbf{How to use it:} Use a delegate for one representative and delegates for more than one.

\textbf{What can follow it:} Use to before the meeting or organization attended, and from before the place or group represented.

\textbf{Common preposition:} Say a delegate to the conference but a delegate from the university.

\textbf{Noun use:} This is a countable noun, so the singular normally needs a, an, or the.

\subsubsection{Useful patterns}
\textbf{Pattern:} a delegate to + conference or meeting

\textbf{Use:} Use this to identify the event or organization that the person attends as a representative.

\textbf{Example:} He was selected as a delegate to the national convention.

\textbf{Pattern:} a delegate from + country, institution, or group

\textbf{Use:} Use this to identify the people or place that the delegate represents.

\textbf{Example:} A delegate from the teachers' union addressed the committee.

\textbf{Pattern:} send or elect a delegate

\textbf{Use:} Use this when a group chooses someone to represent it.

\textbf{Example:} Each department elected a delegate to the university council.

\textbf{Pattern:} conference delegates

\textbf{Use:} Use this for the people officially attending a conference as representatives.

\textbf{Example:} Conference delegates discussed several proposals on the final day.

\subsubsection{Common wrong usage}
\paragraph{Correction 1}
\textbf{Wrong:} She attended the meeting as delegate.

\textbf{Correct:} She attended the meeting as a delegate.

\textbf{Why:} The singular countable noun normally needs a.

\paragraph{Correction 2}
\textbf{Wrong:} He was a delegate of the international conference.

\textbf{Correct:} He was a delegate to the international conference.

\textbf{Why:} Use delegate to for the event attended.

\paragraph{Correction 3}
\textbf{Wrong:} She was a delegate to her university.

\textbf{Correct:} She was a delegate from her university.

\textbf{Why:} Use delegate from for the group or institution represented.

\subsection{Meaning 2: Give someone a task}
\textbf{Part of speech:} transitive verb

\textbf{IPA:} /ˈdɛləˌɡeɪt/

\textbf{Pronunciation phrase:} delegate tasks

\textbf{Local audio file:} audios/delegate-tasks.mp3

\textbf{Register:} neutral to formal; common in business, management, education, and professional contexts

\textbf{Definition:} to give a task or responsibility to another person so that they do it on your behalf

\textbf{In simpler words:} give part of your work to someone else

\textbf{Typical pattern:} delegate + task or responsibility + to + person

\subsubsection{Natural examples}
\begin{enumerate}
  \item Effective managers know when to delegate routine tasks to other team members.
  \item The professor delegated responsibility for collecting the data to two research assistants.
  \item She delegated the preparation of the monthly report to her deputy.
  \item The team leader delegated each stage of the project to a different specialist.
  \item You should delegate some of your work instead of trying to do everything yourself.
  \item The director was reluctant to delegate important responsibilities.
  \item By delegating administrative tasks, doctors can spend more time with patients.
\end{enumerate}

\subsubsection{Nuance and implication}
\textbf{Central nuance:} Delegating means assigning work to another person while usually remaining responsible for the overall result.

\textbf{Usual implication:} Good delegation involves choosing a suitable person and giving them enough information or authority to complete the task.

\textbf{Compared with assign:} Assign is a general word for giving someone a task. Delegate often suggests that the task originally belonged to a person with greater responsibility or authority.

\subsubsection{Grammar notes}
\textbf{How to use it:} This verb normally needs the task or responsibility after it.

\textbf{What can follow it:} It is often followed by task, duty, responsibility, work, decision, or the name of an activity.

\textbf{Common preposition:} Use to before the person receiving the task, as in delegate the work to an assistant.

\textbf{Position in a sentence:} The task normally comes before to and the person: delegate a task to someone.

\subsubsection{Useful patterns}
\textbf{Pattern:} delegate + task or duty + to + person

\textbf{Use:} Use this when giving a particular piece of work to someone else.

\textbf{Example:} The manager delegated the task to an experienced employee.

\textbf{Pattern:} delegate responsibility for + noun or -ing form + to + person

\textbf{Use:} Use this when another person becomes responsible for a particular activity.

\textbf{Example:} She delegated responsibility for organizing the event to her assistant.

\textbf{Pattern:} delegate the preparation or management of + noun

\textbf{Use:} Use this for formal descriptions of work given to another person.

\textbf{Example:} The director delegated the management of the project to a senior engineer.

\textbf{Pattern:} learn or be willing to delegate

\textbf{Use:} Use this without an object when the type of work is already clear.

\textbf{Example:} Successful leaders must learn to delegate.

\subsubsection{Common wrong usage}
\paragraph{Correction 1}
\textbf{Wrong:} The manager delegated to Sara the report.

\textbf{Correct:} The manager delegated the report to Sara.

\textbf{Why:} The usual order is delegate the task to the person.

\paragraph{Correction 2}
\textbf{Wrong:} She delegated her assistant to prepare the report.

\textbf{Correct:} She delegated the preparation of the report to her assistant.

\textbf{Why:} It is usually more natural to delegate a task or responsibility to someone.

\paragraph{Correction 3}
\textbf{Wrong:} He delegated about organizing the conference.

\textbf{Correct:} He delegated responsibility for organizing the conference.

\textbf{Why:} Do not use about after delegate. Name the task or responsibility directly.

\subsection{Meaning 3: Transfer authority}
\textbf{Part of speech:} transitive verb

\textbf{IPA:} /ˈdɛləˌɡeɪt/

\textbf{Pronunciation phrase:} delegate authority

\textbf{Local audio file:} audios/delegate-authority.mp3

\textbf{Register:} formal; common in law, government, administration, and organizational management

\textbf{Definition:} to officially give part of your authority or decision-making power to another person or organization

\textbf{In simpler words:} officially give someone part of your power

\textbf{Typical pattern:} delegate + authority, power, control, or responsibility + to + person or organization

\subsubsection{Natural examples}
\begin{enumerate}
  \item The director delegated authority to approve small purchases to department managers.
  \item National governments sometimes delegate certain powers to regional authorities.
  \item The law allows the agency to delegate some of its responsibilities to local councils.
  \item Senior officials may delegate decision-making authority to local administrators.
  \item The board delegated control of daily operations to the executive director.
  \item Parliament can delegate limited regulatory powers to government agencies.
  \item The president delegated responsibility for the negotiations to the foreign minister.
\end{enumerate}

\subsubsection{Nuance and implication}
\textbf{Central nuance:} This meaning concerns official power rather than merely the performance of a task.

\textbf{Usual implication:} The original authority often keeps overall responsibility and may place limits on the power being transferred.

\textbf{Compared with surrender:} To surrender authority is to give it up, often completely. To delegate authority is usually to transfer only part of it while keeping overall control or responsibility.

\subsubsection{Grammar notes}
\textbf{How to use it:} This verb normally needs the authority, power, control, or responsibility being transferred after it.

\textbf{What can follow it:} Common objects include authority, power, control, responsibility, decision-making, and functions.

\textbf{Common preposition:} Use to before the person, department, or organization receiving the authority.

\textbf{Position in a sentence:} The usual order is delegate authority to someone.

\subsubsection{Useful patterns}
\textbf{Pattern:} delegate authority to + person or organization

\textbf{Use:} Use this when someone officially gives another party the power to make certain decisions.

\textbf{Example:} The minister delegated authority to regional officials.

\textbf{Pattern:} delegate power over + area + to + body

\textbf{Use:} Use this to identify both the area of authority and the organization receiving it.

\textbf{Example:} The law delegates power over local transportation to municipal governments.

\textbf{Pattern:} delegate decision-making responsibility

\textbf{Use:} Use this when another person or group is allowed to make particular choices.

\textbf{Example:} The company delegated decision-making responsibility to local branches.

\textbf{Pattern:} be delegated to + person or organization

\textbf{Use:} Use the passive form when the transferred authority is the main focus.

\textbf{Example:} Certain administrative powers were delegated to the committee.

\subsubsection{Common wrong usage}
\paragraph{Correction 1}
\textbf{Wrong:} The board delegated authority on the manager.

\textbf{Correct:} The board delegated authority to the manager.

\textbf{Why:} Use to for the person or organization receiving authority.

\paragraph{Correction 2}
\textbf{Wrong:} The government delegated the local council with new powers.

\textbf{Correct:} The government delegated new powers to the local council.

\textbf{Why:} The usual pattern is delegate something to someone.

\paragraph{Correction 3}
\textbf{Wrong:} The director delegated all responsibility and was no longer accountable.

\textbf{Correct:} The director delegated several responsibilities but remained accountable for the final results.

\textbf{Why:} Delegating work or authority does not necessarily remove the original person's overall responsibility.

\section{Word family}
\sectionrule
\subsection{delegation}
\textbf{Part of speech:} countable or uncountable noun

\textbf{IPA:} /ˌdɛləˈɡeɪʃən/

\textbf{Definition:} a group of people chosen to represent others, or the act of giving tasks or authority to another person

\textbf{Relationship to the headword:} This noun can refer either to a group of delegates or to the process of delegating work or authority.

\textbf{Grammar pattern:} a delegation from + place, a delegation of + representatives, or the delegation of + task or authority

\textbf{Usage note:} Use a delegation for a representative group. Delegation is often uncountable when it means the management process.

\subsubsection{Examples}
\begin{enumerate}
  \item A delegation of scientists met with government officials.
  \item Effective delegation allows managers to focus on their most important responsibilities.
  \item The international delegation visited several research institutions.
  \item Poor delegation can leave employees confused about their responsibilities.
\end{enumerate}

\section{Common collocations}
\sectionrule
\subsection{conference delegate}
\textbf{Applies to:} Chosen representative

\textbf{Meaning:} a person officially attending a conference as a representative

\textbf{Example:} Each conference delegate received a copy of the research summary.

\subsection{delegate to a conference}
\textbf{Applies to:} Chosen representative

\textbf{Meaning:} a representative chosen to attend a conference

\textbf{Pattern:} a delegate to + meeting, conference, or convention

\textbf{Example:} She was selected as a delegate to an international education conference.

\subsection{delegate from a country}
\textbf{Applies to:} Chosen representative

\textbf{Meaning:} a representative of a particular country or group

\textbf{Pattern:} a delegate from + country, institution, or organization

\textbf{Example:} A delegate from Brazil presented the proposal.

\subsection{delegate a task}
\textbf{Applies to:} Give someone a task

\textbf{Meaning:} give a piece of work to another person

\textbf{Pattern:} delegate a task to + person

\textbf{Example:} The supervisor delegated the task to a junior researcher.

\subsection{delegate responsibility}
\textbf{Applies to:} Give someone a task

\textbf{Meaning:} make another person responsible for completing particular work

\textbf{Pattern:} delegate responsibility for + noun or -ing form + to + person

\textbf{Example:} Managers should delegate responsibility without losing oversight.

\subsection{delegate effectively}
\textbf{Applies to:} Give someone a task

\textbf{Meaning:} give tasks to suitable people in a clear and useful way

\textbf{Example:} Leaders who delegate effectively can use their time more efficiently.

\subsection{delegate authority}
\textbf{Applies to:} Transfer authority

\textbf{Meaning:} officially transfer some decision-making power

\textbf{Pattern:} delegate authority to + person or organization

\textbf{Example:} The chief executive delegated authority to several regional managers.

\subsection{delegate power}
\textbf{Applies to:} Transfer authority

\textbf{Meaning:} officially give another person or organization certain powers

\textbf{Pattern:} delegate power to + person or organization

\textbf{Example:} The constitution permits the national government to delegate limited powers to local authorities.

\subsection{delegated authority}
\textbf{Applies to:} Transfer authority

\textbf{Meaning:} authority officially transferred from a higher authority

\textbf{Example:} The officer acted within the limits of her delegated authority.

\textbf{Usage note:} Delegated is an adjective in this expression.

\subsection{a delegation of officials}
\textbf{Applies to:} Chosen representative

\textbf{Meaning:} a representative group of officials

\textbf{Pattern:} a delegation of + representatives

\textbf{Example:} A delegation of officials traveled to the neighboring country.

\subsection{effective delegation}
\textbf{Applies to:} Give someone a task

\textbf{Meaning:} the successful assignment of work and responsibility to other people

\textbf{Example:} Effective delegation improves both productivity and staff development.

\textbf{Usage note:} Delegation is the noun form.

\section{Synonyms and antonyms}
\sectionrule
\subsection{Close synonyms}
\subsubsection{representative}
\textbf{Part of speech:} countable noun

\textbf{Applies to:} Chosen representative

\textbf{Shared meaning:} Both words describe someone who acts on behalf of a larger group.

\textbf{Difference:} Representative is a broad term for anyone who acts or speaks for others. Delegate often refers specifically to someone formally chosen to attend a meeting, conference, or organization.

\textbf{Example:} Each organization sent a representative to the meeting.

\subsubsection{assign}
\textbf{Part of speech:} transitive verb

\textbf{Applies to:} Give someone a task

\textbf{Shared meaning:} Both verbs can mean giving work to another person.

\textbf{Difference:} Assign simply means give someone a task. Delegate often suggests that a person passes on work that originally belonged to them while keeping overall responsibility.

\textbf{Example:} The professor assigned each student a separate research topic.

\subsubsection{entrust}
\textbf{Part of speech:} transitive verb

\textbf{Applies to:} Give someone a task

\textbf{Shared meaning:} Both verbs can describe giving another person responsibility.

\textbf{Difference:} Entrust emphasizes confidence that someone will take care of an important task, object, or person. Delegate focuses on transferring responsibility for work.

\textbf{Example:} The director entrusted her assistant with the confidential documents.

\subsubsection{authorize}
\textbf{Part of speech:} transitive verb

\textbf{Applies to:} Transfer authority

\textbf{Shared meaning:} Both verbs can involve officially giving someone power to act.

\textbf{Difference:} Authorize means officially give permission to do something. Delegate means transfer some of your own authority or responsibility to another party.

\textbf{Example:} The committee authorized the director to sign the agreement.

\section{Words learners may confuse}
\sectionrule
\subsection{designate}
\textbf{Part of speech:} transitive verb

\textbf{Why learners confuse them:} Both formal verbs can involve giving someone a role, but they describe different actions.

\textbf{Key difference:} Delegate means give someone a task or authority. Designate usually means officially choose, name, or identify someone or something for a role or purpose.

\textbf{delegate:} The director delegated responsibility for the project to Maria.

\textbf{designate:} The director designated Maria as the project leader.

\section{Etymology}
\sectionrule
\textbf{Origin:} Delegate comes from Latin delegare, meaning to send as a representative or assign responsibility.

\section{Practice exercises}
\sectionrule
\subsection{Word family choice}
Choose the best member of the word family for each blank.

\begin{enumerate}
  \item A \_\_\_\_\_ of university officials attended the international meeting.
  \item Effective \_\_\_\_\_ can prevent managers from becoming overwhelmed.
  \item Each department selected a \_\_\_\_\_ to attend the council meeting.
\end{enumerate}

\subsection{Correct or incorrect?}
Decide whether each sentence is correct. Correct the inaccurate ones.

\begin{enumerate}
  \item The manager delegated the task to an experienced employee.
  \item She attended the conference as delegate from her university.
  \item The board delegated authority on the regional director.
  \item Good leaders are willing to delegate.
\end{enumerate}

\subsection{Collocation practice}
Complete each sentence with a natural collocation.

\begin{enumerate}
  \item The supervisor decided to delegate the task \_\_\_\_\_ a more experienced employee.
  \item The two departments each elected a \_\_\_\_\_ to the university council.
  \item The law allows the ministry to delegate \_\_\_\_\_ to local agencies.
  \item Effective \_\_\_\_\_ requires clear instructions and appropriate supervision.
\end{enumerate}

\subsection{Complete answer key}
\subsubsection{Word family choice}
\begin{enumerate}
  \item \textbf{delegation.} A delegation of university officials attended the international meeting. \emph{Use delegation for a representative group.}
  \item \textbf{delegation.} Effective delegation can prevent managers from becoming overwhelmed. \emph{Use delegation for the process of giving work to other people.}
  \item \textbf{delegate.} Each department selected a delegate to attend the council meeting. \emph{Use delegate for the chosen representative.}
\end{enumerate}

\subsubsection{Correct or incorrect?}
\begin{enumerate}
  \item \textbf{Correct} \emph{Delegate a task to someone is the standard pattern.}
  \item \textbf{Incorrect}. She attended the conference as a delegate from her university. \emph{The singular countable noun requires a.}
  \item \textbf{Incorrect}. The board delegated authority to the regional director. \emph{Use to before the person receiving the authority.}
  \item \textbf{Correct} \emph{Delegate can appear without an object when the type of work is clear from the context.}
\end{enumerate}

\subsubsection{Collocation practice}
\begin{enumerate}
  \item \textbf{to.} The supervisor decided to delegate the task to a more experienced employee. \emph{Use delegate something to someone.}
  \item \textbf{delegate.} The two departments each elected a delegate to the university council. \emph{A delegate to a council is a person chosen to attend it as a representative.}
  \item \textbf{authority.} The law allows the ministry to delegate authority to local agencies. \emph{Delegate authority means officially transfer some decision-making power.}
  \item \textbf{delegation.} Effective delegation requires clear instructions and appropriate supervision. \emph{Effective delegation is the successful assignment of work to others.}
\end{enumerate}

\vfill
\begin{center}
\small\color{Muted} Created and compiled by \textbf{Amir Kooshky}\\
\href{https://t.me/KooshkyTOEFL}{Telegram Channel}\quad\textbullet\quad
\href{https://www.instagram.com/kooshkytoefl}{Instagram}\quad\textbullet\quad
\href{https://t.me/KooshkyTOEFL\_pv}{Personal Telegram}
\end{center}
\end{document}
